\section{Stability of knot equivalence}\label{sec:stability}
The equivalence of two tame knots $\g_1,\g_2\in C^0(\R/\Z,\R^3)$
with local distortion below $g_3$ (or even below
$\frac{\pi}{\sqrt{8}}$) at some positive scales, such that the knots 
are sufficiently close in 
terms of these scales, was stated in Corollary \ref{cor:equiv-knots}. This was a direct
consequence of Theorem~\ref{thm:isomorphic} applied to the knots' images $M_i:=
\g_i(\R/\Z)$, $i=1,2,$ in combination with Theorem~\ref{thm:gordon-luecke}. 
We are going to apply Corollary \ref{cor:equiv-knots} now
to a knot and a small Lipschitz perturbation of that knot to
prove the Lipschitz stability of knot equivalence, i.e., 
Corollary \ref{cor:lipschitz-stability}. For that we specify the local distortion
at some scale $r>0$ as
defined in \eqref{eq:local-distortion} in Definition \ref{def:local-distortion}
to the image $M:=\eta(\R/\Z)$ of a knot $\eta\in C^0(\R/\Z,\R^n)$ by writing
\begin{equation}\label{eq:local-distortion-knots}
\delta(\eta,r)\equiv\delta(\eta(\R/\Z),r) 
=\sup_{u,v\in\R/\Z, u\not= v\atop 0<|\eta(u)-\eta(v)|\le 2r}
\frac{d_\eta\big(\eta(u),\eta(v)\big)}{|\eta(u)-\eta(v)|},
\end{equation}
where the intrinsic distance $d_\eta(\eta(u),\eta(v))$ reduces to the periodic 
parameter distance $|u-v|_{\R/\Z}$ in case of arclength parametrized curves $\eta$.

\begin{proof}[Proof of Corollary {\ref{cor:lipschitz-stability}}]
By assumption there is a scale $r_\g\in(0,\tfrac12)$ such that $\delta(\g,r_\g)< g_3$ 
(or even
less than $\frac{\pi}{\sqrt{ 8 }}$.) Our aim is to show that for a given
function $h\in C^{0,1}(\R/\Z,\R^3)$ the perturbed
curve $\g_\epsilon=\g+\epsilon h$ has also local distortion
$\delta(\g_\epsilon,\hat{r}_{\g})<g_3 $ at
a (possibly smaller) scale $\hat{r}_\g>0$ independent of $0<|\epsilon|\ll 1$. Keeping
that scale fixed we can then take $|\epsilon| $ so small 
(depending on~$h$) that the $L^\infty$-distance
is controlled in terms of the two scales $r_\g$ and $\hat{r}_\g$.

To that extent consider $x,y\in [0,1]$ with $x<y$ 
to estimate
\begin{equation}\label{eq:euclid-eps}
|\g_\epsilon(x)-\g_\epsilon(y)|=
|\g(x)-\g(y)+\epsilon(h(x)-h(y))|\ge |\g(x)-\g(y)|-|\epsilon|\|h'\|_{L^\infty}(y-x).
\end{equation}
The intrinsic distance $d_{\g_\epsilon}(\g_\epsilon(x),\g_\epsilon(y))\equiv
d_{\g_\epsilon}(x,y)$ on $\g_\epsilon$ may be compared
to $d_\g(\g(x),\g(y))\equiv d_\g(x,y)$ on $\g$ by means of
\begin{align}\label{eq:intrinsic-eps-prep}
\textstyle
\big|\int_x^y |\g_\epsilon'(u)|du & -
\textstyle
\int_x^y |\g'(u)|du\big|=
\big|\int_x^y \big(|\g_\epsilon'(u)|-|\g'(u)|\big)du\big|\le\int_x^y
\big| |\g_\epsilon'(u)|-|\g'(u)|\big|du\notag\\
&\textstyle
\le\int_x^y |\g_\epsilon'(u)-\g'(u)|du=
\int_x^y|\epsilon h'(u)|du\le |\epsilon|\|h'\|_{L^\infty}(y-x),
\end{align}
% so that in case $0<y-x\le\frac13$ and with $\epsilon_1:=(2\|h'\|_{L^\infty}+1)^{-1}$
% $$
% \textstyle
% \frac13\ge y-x=\int_x^y |\g'(u)|du\overset{\eqref{eq:intrinsic-eps-prep}}{\ge}
% \int_x^y |\g_\epsilon'(u)|du -\frac{|\epsilon|}3 \|h'\|_{L^\infty}
% >\int_x^y |\g_\epsilon'(u)|du-\frac16\Foa \epsilon\in [-\epsilon_1,\epsilon_1];
% $$
whence
\begin{equation}\label{eq:intrinsic-eps}
\textstyle
d_{\g_\epsilon}(x,y)\le\int_x^y|\g_\epsilon'(u)|du
\overset{\eqref{eq:intrinsic-eps-prep}}{\le}
\int_x^y|\g'(u)|du+|\epsilon|\|h'\|_{L^\infty}(y-x)
\le(y-x)(1+|\epsilon|\|h'\|_{L^\infty}).
\end{equation}
Without loss of generality we may assume that 
$0< y-x\le\tfrac12$, so
$d_\g(x,y) = y-x$.
Combining \eqref{eq:euclid-eps} with \eqref{eq:intrinsic-eps} we deduce supposing
 $|\g(x)-\g(y)|<2r_\g$
\begin{align}
\begin{split}\label{eq:distortion-eps}
\textstyle\frac{d_{\g_\epsilon}(x,y)}{|\g_\epsilon(x)-\g_\epsilon(y)|}
&  \overset{\eqref{eq:intrinsic-eps},\eqref{eq:euclid-eps}}{\le}
\textstyle
\frac{d_{\g}(x,y)(1+|\epsilon|\|h'\|_{L^\infty})}{|\g(x)-\g(y)|-
|\epsilon|\|h'\|_{L^\infty}(y-x)}=
 \frac{d_{\g}(x,y)(1+|\epsilon|\|h'\|_{L^\infty})}{|\g(x)-\g(y)|
 \big(1-|\epsilon|\|h'\|_{L^\infty}\frac{d_\g(x,y)}{|\g(x)-\g(y)|}\big)} \\
 & \le \textstyle
 \delta(\g,r_\g)\cdot\frac{1+|\epsilon|\|h'\|_{L^\infty}}{
 1-|\epsilon|\|h'\|_{L^\infty}\delta(\g,r_\g)}<g_3,
\end{split}
\end{align}
if
\begin{equation}\label{eq:eps2}
\textstyle
|\epsilon|\le\epsilon_0:=\min\big\{\frac{g_3-\delta(\g,r_\g)}{\delta(\g,r_\g)(1+g_3)\|h'\|_{L^\infty}+1},\frac{r_\g}{\|h'\|_{L^\infty}+1}\big\}>0.
% \min\big\{
% \frac{g_3-\delta(\g,r_\g)}{g_3\|h'\|_{L^\infty}(1+g_3)+1},\frac1{g_3\|h'\|_{L^\infty}+1},
% \epsilon_1\big\}.
\end{equation}
% For the second inequality in \eqref{eq:distortion-eps} we have also used the fact
% that the
% function $s\mapsto s(1+c)/(1-cs)$ is non-decreasing on $[0,\frac1{c})$ for any
% constant $c\in (0,1)$.
Consequently, if $|\g_\epsilon(x)-\g_\epsilon(y)|\le2\hat r_\g$ where $\hat r_\g=\tfrac12 r_\g$
we infer
$|\g(y)-\g(x)|\le 2r_\g$, so $\delta(\g_\epsilon,\hat{r}_\g)<g_3$
for all $|\epsilon|\le\epsilon_0$
due to \eqref{eq:distortion-eps} and
\eqref{eq:eps2}.
% For $y-x\in [\Delta_0,\frac12]$ we use that
% $\sigma_\g:=\min\big\{|\g(y)-\g(x)|:\Delta_0\le d_\g(x,y)\big\}\in (0,\frac13]$,
% since $\g$ is embedded. For $|\epsilon|\le\epsilon_3:=\frac{\sigma_\g}{\|h'\|_{L^\infty}+1}$ we obtain from \eqref{eq:euclid-eps}
% \begin{equation}\label{eq:euclid-eps-global}
% \textstyle
% |\g_\epsilon(x)-\g_\epsilon(y)|\overset{\eqref{eq:euclid-eps}}{\ge}
% \sigma_\g-|\epsilon|\|h'\|_{L^\infty}\frac12 >
% \frac{\sigma_\g}2\Fo y-x\in [\Delta_0,\frac12].
% \end{equation}
% Combining all requirements on $|\epsilon|$ so far we set $\epsilon_4:=\min\{
% \epsilon_2,\epsilon_3\}$ and $\hat{r}_\g:=\frac{\sigma_\g}4$, to obtain
% $\delta(\g_\epsilon,\hat{r}_\g)<g_3$  
% for all $|\epsilon|\le\epsilon_4$ due to \eqref{eq:distortion-eps} and
% \eqref{eq:euclid-eps-global}.
It is important to notice that the scale
$\hat{r}_\g$ does \emph{not} depend on $\epsilon\in [-\epsilon_0,\epsilon_0]$.
With $\|\g-\g_\epsilon\|_{L^\infty}=|\epsilon|\|h\|_{L^\infty}$ this allows
us to choose
\begin{equation}\label{eq:final-epsilon-choice}
\textstyle
\epsilon_{\g,h}:=\min\Big\{\epsilon_0,\frac{\hat{r}_\g}{8\|h \|_{
L^\infty}+1}\Big\}
<\frac14 \min\{r_\g,\hat{r}_\g\}
\end{equation}
to also satisfy the hypothesis on the $L^\infty$-distance in Corollary 
\ref{cor:equiv-knots} to conclude that $\g\sim\g_\epsilon$ for all
$|\epsilon|\le\epsilon_{\g,h}.$ In addition, $\essinf_{\R/\Z}|\g_\epsilon'|
\ge\essinf_{\R/\Z}|\g'|-|\epsilon|\|h'\|_{L^\infty} > \frac12$
for all $|\epsilon|\le\epsilon_{\g,h}$, since $\epsilon_{\g,h}\le\epsilon_0\le
\tfrac12(\|h'\|_{L^\infty}+1)^{-1}$.
\end{proof}

% \noindent
% {\it Proof of Corollary \ref{cor:lipschitz-stability}.}\,
% By assumption there is a scale $r_\g>0$ such that $\delta(\g,r_\g)< g_3$ 
% (or even
% less than $\frac{\pi}{\sqrt{ 8 }}$.) Our aim is to show that for a given
% function $h\in C^{0,1}(\R/\Z,\R^3)$ the perturbed
% curve $\g_\epsilon=\g+\epsilon h$ has also local distortion
% $\delta(\g_\epsilon,\hat{r}_{\g})<g_3 $ at
% a (possibly smaller) scale $\hat{r}_\g>0$ independent of $0<|\epsilon|\ll 1$. Keeping
% that scale fixed we can then take $|\epsilon| $ so small \philipp{(depending on~$h$)} that the $L^\infty$-distance
% is controlled in terms of the two scales $r_\g$ and $\hat{r}_\g$.

% To that extent consider $x,y\in [0,1]$ with $0<y-x\le 1$ to estimate
% \begin{equation}\label{eq:euclid-eps}
% |\g_\epsilon(x)-\g_\epsilon(y)|=
% |\g(x)-\g(y)+\epsilon(h(x)-h(y))|\ge |\g(x)-\g(y)|-|\epsilon|\|h'\|_{L^\infty}(y-x).
% \end{equation}
% The intrinsic distance $d_{\g_\epsilon}(\g_\epsilon(x),\g_\epsilon(y))\equiv
% d_{\g_\epsilon}(x,y)$ on $\g_\epsilon$ may be compared
% to $d_\g(\g(x),\g(y))\equiv d_\g(x,y)$ on $\g$ by means of
% \begin{align}\label{eq:intrinsic-eps-prep}
% \textstyle
% \big|\int_x^y |\g_\epsilon'(u)|du & -
% \textstyle
% \int_x^y |\g'(u)|du\big|=
% \big|\int_x^y \big(|\g_\epsilon'(u)|-|\g'(u)|\big)du\big|\le\int_x^y
% \big| |\g_\epsilon'(u)|-|\g'(u)|\big|du\notag\\
% &\textstyle
% \le\int_x^y |\g_\epsilon'(u)-\g'(u)|du=
% \int_x^y|\epsilon h'(u)|du\le |\epsilon|\|h'\|_{L^\infty}(y-x),
% \end{align}
% so that in case $0<y-x\le\frac13$ and with $\epsilon_1:=(2\|h'\|_{L^\infty}+1)^{-1}$
% $$
% \textstyle
% \frac13\ge y-x=\int_x^y |\g'(u)|du\overset{\eqref{eq:intrinsic-eps-prep}}{\ge}
% \int_x^y |\g_\epsilon'(u)|du -\frac{|\epsilon|}3 \|h'\|_{L^\infty}
% >\int_x^y |\g_\epsilon'(u)|du-\frac16\Foa \epsilon\in [-\epsilon_1,\epsilon_1];
% $$
% whence\phtodo{den ganzen Beweis nochmal checken}
% \begin{equation}\label{eq:intrinsic-eps}
% \textstyle
% d_{\g_\epsilon}(x,y)\le\int_x^y|\g_\epsilon'(u)|du
% \overset{\eqref{eq:intrinsic-eps-prep}}{\le}
% \int_x^y|\g'(u)|du+|\epsilon|\|h'\|_{L^\infty}(y-x)
% =d_\g(x,y)+|\epsilon|\|h'\|_{L^\infty}(y-x).
% \end{equation}
% Combining \eqref{eq:euclid-eps} with \eqref{eq:intrinsic-eps} we deduce
% for $0<y-x<\Delta_0:=\min\{2r_\g,\frac13\}$ (so that $|\g(x)-\g(y)|<2r_\g$)
% \begin{eqnarray}\label{eq:distortion-eps}
% \textstyle\frac{d_{\g_\epsilon}(x,y)}{|\g_\epsilon(x)-\g_\epsilon(y)|}
% &  \overset{\eqref{eq:intrinsic-eps},\eqref{eq:euclid-eps}}{\le}  & 
% \textstyle
% \frac{d_{\g}(x,y)+|\epsilon|\|h'\|_{L^\infty}(y-x)}{|\g(x)-\g(y)|-
% |\epsilon|\|h'\|_{L^\infty}(y-x)}=
%  \frac{d_{\g}(x,y)\big(1+|\epsilon|\|h'\|_{L^\infty}\big)}{|\g(x)-\g(y)|
%  \big(1-|\epsilon|\|h'\|_{L^\infty}\frac{d_\g(x,y)}{|\g(x)-\g(y)|}\big)}\notag\\
%  & \le & \textstyle
%  \delta(\g,r_\g)\cdot\frac{1+|\epsilon|\|h'\|_{L^\infty}}{
%  1-|\epsilon|\|h'\|_{L^\infty}\delta(\g,r_\g)}<g_3,
% \end{eqnarray}
% if
% \begin{equation}\label{eq:eps2}
% \textstyle
% |\epsilon|\le\epsilon_2:=\min\big\{
% \frac{g_3-\delta(\g,r_\g)}{g_3\|h'\|_{L^\infty}(1+g_3)+1},\frac1{g_3\|h'\|_{L^\infty}+1},
% \epsilon_1\big\}.
% \end{equation}
% For the second inequality in \eqref{eq:distortion-eps} we have also used the fact
% that the
% function $s\mapsto s(1+c)/(1-cs)$ is non-decreasing on $[0,\frac1{c})$ for any
% constant $c\in (0,1)$.

% For $y-x\in [\Delta_0,\frac12]$ we use that
% $\sigma_\g:=\min\big\{|\g(y)-\g(x)|:\Delta_0\le d_\g(x,y)\big\}\in (0,\frac13]$,
% since $\g$ is embedded. For $|\epsilon|\le\epsilon_3:=\frac{\sigma_\g}{\|h'\|_{L^\infty}+1}$ we obtain from \eqref{eq:euclid-eps}
% \begin{equation}\label{eq:euclid-eps-global}
% \textstyle
% |\g_\epsilon(x)-\g_\epsilon(y)|\overset{\eqref{eq:euclid-eps}}{\ge}
% \sigma_\g-|\epsilon|\|h'\|_{L^\infty}\frac12 >
% \frac{\sigma_\g}2\Fo y-x\in [\Delta_0,\frac12].
% \end{equation}
% Combining all requirements on $|\epsilon|$ so far we set $\epsilon_4:=\min\{
% \epsilon_2,\epsilon_3\}$ and $\hat{r}_\g:=\frac{\sigma_\g}4$, to obtain
% $\delta(\g_\epsilon,\hat{r}_\g)<g_3$  
% for all $|\epsilon|\le\epsilon_4$ due to \eqref{eq:distortion-eps} and
% \eqref{eq:euclid-eps-global}. It is important to notice that the scale
% $\hat{r}_\g$ does \emph{not} depend on $\epsilon\in [-\epsilon_4,\epsilon_4]$.
% With $\|\g-\g_\epsilon\|_{L^\infty}=|\epsilon|\|h\|_{L^\infty}$ this allows
% us to choose
% \begin{equation}\label{eq:final-epsilon-choice}
% \textstyle
% \epsilon_{\g,h}:=\min\Big\{\epsilon_4,\frac14 \min\{r_\g,\hat{r}_\g\}\cdot\frac1{\|h \|_{
% L^\infty}+1}\Big\}
% \end{equation}
% to also satisfy the hypothesis on the $L^\infty$-distance in Corollary 
% \ref{cor:equiv-knots} to conclude that $\g\sim\g_\epsilon$ for all
% $|\epsilon|\le\epsilon_{\g,h}.$ In addition, $\essinf_{\R/\Z}|\g_\epsilon'|
% \ge\essinf_{\R/\Z}|\g'|-|\epsilon|\|h'\|_{L^\infty} > \frac12$
% for all $|\epsilon|\le\epsilon_{\g,h}$, since $\epsilon_{\g,h}\le\epsilon_1=
% (2\|h'\|_{L^\infty}+1)^{-1}.$ \hfill $\Box$

Before treating the stability of knot equivalence in the setting of
fractional Sobolev regularity let us point out that bounding  the
bilipschitz constant\footnote{\label{foot:bilip}In \cite{freedman-etal_1994} a bilipschitz constant
is defined as any number $L>0$ such that $L^{-1}d_\gamma(u,v)|\le|\gamma(u)-\gamma(v)|\le Ld_\gamma(u,v)$
holds for all $u,v\in\R/\Z$.
If $\gamma$ is parametrized by arclength then the infimum of those $L$
is $\bilip_I(\g)^{-1}$.}
\begin{equation}\label{eq:def-bilip}
\bilip_I(\g):=\inf_{\substack{u,v\in I\\u\ne v}}\frac{|\g(u)-\g(v)|}{|u-v|_{\R/\Z}}
\end{equation}
on an interval $I\subset\R$ from below, is closely related to
controlling the local distortion $\delta(\g,\cdot)$ at some
scale related to $I$, if the absolutely continuous parametrization $\g$
has positive
speed $v_\g=\essinf_{\R/\Z}|\g'|$, i.e., if $\g$ is an immersed curve.

The seminorm $\lfloor \g'\rfloor $ defining fractional Sobolev spaces allows us to bound
$\bilip(\g)$ as  in the following lemma, which is essentially contained
in \cite[Lemma 2.1]{blatt_2012a}; see also \cite[Lemma 2.1]{blatt-etal_2024}.


\begin{lemma}\label{lem:seminorm-bilip}
Let $I\subset\R$ be an interval and $\g:I\to\R^n$ an absolutely continuous
immersed curve with positive speed $v_\g=\essinf_{I}|\g'|$. 
Then one has for all $s,t\in I$ with $s<t$
\begin{equation}\label{eq:seminorm-bilip}
|\g(t)-\g(s)|^2\ge\Big\{v_\g^2-\frac12\int_s^t\int_s^t\frac{|\g'(u)-\g'(v)|^2}{|u-v|^2}dudv\Big\}|t-s|^2.
\end{equation}
\end{lemma}
As an immediate consequence one deduces
\begin{corollary}[Fractional seminorm controls bilipschitz constant]
\label{cor:seminorm-bilip}
If $\g\in W^{\frac32, 2}(I,\R^n)$ is arclength parametrized, then
\begin{equation}\label{eq:seminorm-bilip2}
|\g(t)-\g(s)|^2\ge\big(1-\frac12 \lfloor\g'\rfloor^2_{\frac12,2,I}\big)|t-s|^2\quad\Foa
s,t\in I.
\end{equation}
\end{corollary}
%\noindent
\begin{proof}[Proof of Lemma {\ref{lem:seminorm-bilip}}]
%{\it Proof of Lemma \ref{lem:seminorm-bilip}.}\,
For parameters $s,t\in I$ with $s<t$ estimate
\begin{gather*}
|\g(s) -\g(t)|^2  
\textstyle
= \int_s^t\int_s^t\langle\g'(\sigma),\g'(\tau)\rangle_{\R^n}d\sigma 
d\tau
 \textstyle
= \frac12  \int_s^t\int_s^t \big(
|\g'(\tau)|^2+
|\g'(\sigma)|^2-|\g'(\tau)-\g'(\sigma)|^2\big)d\sigma d\tau\\
\textstyle
\ge v_\g^2|t-s|^2-\frac12\int_s^t\int_s^t|\g'(\tau)-\g'(\sigma)|^2d\sigma d\tau
\ge\textstyle
|t-s|^2\big\{v_\g^2-\frac12\int_s^t\int_s^t
\frac{|\g'(\tau)-\g'(\sigma)|^2}{|\tau-\sigma|^2}d\sigma d\tau\big\}.\qedhere
\end{gather*}
%\hfill $\Box$
\end{proof}

We observe in \eqref{eq:seminorm-bilip2} that a small fractional seminorm
yields a bilipschitz  estimate, which is the key ingredient for the proof 
of the fractional stability of knot equivalence.

\noindent
{\it Proof of Corollary \ref{cor:fractional-stability}.}\,
%For the dimension-dependent constant $g_n$ from  \eqref{eq:dekster-jung}, 
%in 
%particular for $g_3=\frac{\sqrt{32}}{\sqrt{27}}\frac{\pi}{\sqrt{8}}\in 
%(\frac2{\sqrt{3}},\frac{\pi}2)$ in 
%Remark \ref{rem:dimension-constant}, set
%$$
%c_3:=\sqrt{2-\frac2{g_3^2}}=\sqrt{2-\frac{27}{2\pi^2}}\in \Big(\sqrt{\frac12},1\Big).
%$$
Due to the absolute continuity of the double integral defining the
fractional seminorm there is a radius $\rho=\rho(\g)\in (0,\frac14]$  such
that $\sup_{x\in\R/\Z}\lfloor \g'\rfloor_{\frac12,2,B_\rho(x)} <\frac1{2\sqrt{2}}.$
Therefore,
\begin{equation}\label{eq:eta-seminorm}
\textstyle
\sup_{x\in\R/\Z}\lfloor \eta' \rfloor_{\frac12,2,B_\rho(x)}
\le
\sup_{x\in\R/\Z}\lfloor \g' \rfloor_{\frac12,2,B_\rho(x)}+\lfloor
\eta'-\g' \rfloor_{\frac12,2} < \frac1{\sqrt{2}}
\end{equation}
for all $\eta \in B_{\frac1{2\sqrt{2}}}(\g)\subset W^{\frac32,2}(\R/\Z,\R^3)$. By
means of Corollary \ref{cor:seminorm-bilip} we infer
\begin{equation}\label{eq:bilip-gamma-eta}
\textstyle
\min\big\{|\g(t)-\g(s)|,|\eta(t)-\eta(s)|\big\}>\frac{\sqrt{3}}2 
|t-s|\ge\frac{\sqrt{3}}2 d_\g(t,s)
=\frac{\sqrt{3}}2 d_\eta(t,s)
\end{equation}
for all arclength parametrized curves $\eta\in B_{\frac1{2\sqrt{2}}}(\g)\subset
W^{\frac32,2}(\R/\Z,\R^3)$ and all parameters $s,t\in\R/\Z$ with
intrinsic distance $d_\g(t,s)\equiv d_\g(\g(s),\g(t))=|t-s|_{\R/\Z}=
d_\eta(\eta(s),\eta(t))\equiv d_\eta(t,s)<2\rho.$

On the other hand,
$$
\textstyle
\sigma_\g:=
\inf\{|\g(t)-\g(s)|:d_\g(s,t)=|t-s|_{\R/\Z}\ge 2\rho\}\in (0,\frac12].
$$
By Sobolev embedding there is a constant $C_E$ such that $\|\zeta\|_{L^\infty}
\le C_E\|\zeta\|_{W^{\frac32,2}}$ for all $\zeta\in W^{\frac32,2}(\R/\Z,\R^3)$,
so that
\begin{equation}\label{eq:infty-distance-eta-gamma}
\textstyle
\|\eta-\gamma\|_{L^\infty}< \frac{\sigma_\g}{16}\quad\Foa\eta\in B_{\epsilon_\g}(\g)
\subset W^{\frac32,2}(\R/\Z,\R^3),
\end{equation}
where we have set $\epsilon_\g:=\min\{\frac1{2\sqrt{2}},\frac{\sigma_\g}{16C_E}\}$.
Therefore, for all $s,t\in\R/\Z$ with $d_\g(s,t)=d_\eta(s,t)
\ge 2\rho$ one has
\begin{equation}\label{eq:euclid-eta-gamma-global}
\textstyle
\min\big\{|\g(t)-\g(s)|,|\eta(t)-\eta(s)|\big\}>\frac{7}{8}\sigma_\g
>\frac{1}{2}\sigma_\g.
\end{equation}
Defining the common scale $r_{\g}:=\frac{\sigma_\g}4$ we find in view of 
\eqref{eq:bilip-gamma-eta} and \eqref{eq:euclid-eta-gamma-global}
\begin{equation}\label{eq:max-eta-gamma-distortion}
\textstyle
\max\big\{\delta(\g,r_\g),\delta(\eta,r_\g)\big\}\le\frac{2}{\sqrt{3}}<g_3
\end{equation}
for all arclength parametrized curves $\eta\in B_{\epsilon_\g}(\g)\subset
W^{\frac32,2}(\R/\Z,\R^3)$, since 
$g_3=\frac{\sqrt{32}}{\sqrt{27}}\frac{\pi}{\sqrt{8}}\in 
(\frac2{\sqrt{3}},\frac{\pi}2)$.
Inequality \eqref{eq:infty-distance-eta-gamma} allows us to apply
Corollary \ref{cor:equiv-knots} to conclude. Notice that we also used
the fact that arclength parametrized $W^{\frac32,2}$-knots have finite
M\"obius energy according to \cite[Theorem~1.1]{blatt_2012a}, and
are therefore tame \cite[Theorem~4.1]{freedman-etal_1994}, which is needed
in Corollary \ref{cor:equiv-knots}.\hfill $\Box$
